\chapter{SYSTEM ANALYSIS AND FEASIBILITY}
System analysis helps bridge design and problem identification. It tells us what must be done, the limitations and constraints the system will be bound within, and the requirements for measuring the quality of the outputs the system produces. In this chapter, project goals are described, the methods and processes of requirements elicitation are detailed, and the proposed solution is examined for feasibility in terms of technology, law, economy, and labor.

\section{System Analysis}
As this is an R&D project, requirements were not provided because there were no client specifications. Instead, requirements were developed through the analysis of issues and constraints provided by existing deep learning based HSI classifier algorithms. The main function of the system is to read and handle high-dimensional spectral data (either introduced externally by the user or extracted from established benchmarks - Indian Pines, Pavia University, and Salinas) and classify high accuracy pixel data.

The system is governed by three functional requirements:
\begin{itemize}
    \item Raw hyperspectral cubes must be filtered using a Gaussian filter and PCA \cite{pedregosa2011scikit}.
    \item Supervised (CNN \cite{chen2016deep}, ViT \cite{dosovitskiy2020image}) and semi-supervised (Self-Training \cite{triguero2015self}) models must be trained in a Curriculum Learning framework \cite{bengio2009curriculum} that regulates the difficulty of samples.
    \item A multithreaded GUI must be used to encapsulate all trained models, allowing users to perform simultaneous model testing and view per pixel predictions in real-time without modifying the underlying source code.
\end{itemize}

Meeting these requirements will be evaluated using Accuracy, Precision, Recall, F1-Score, and Cohen’s Kappa coefficient.

\section{Feasibility}
The analysis examines the likelihood of the proposed solution in four categories: technical, legal, economic, and workforce.

\subsection{Technical Feasibility}
The proposed programming language is Python \cite{van1995python} since it has a complete ecosystem for machine learning with a favorable community. Deep learning is constructed with PyTorch \cite{paszke2019pytorch} or TensorFlow \cite{abadi2016tensorflow} with data management offered by Scikit-learn \cite{pedregosa2011scikit} and OpenCV \cite{itseez2015opencv}. The design of the graphical user interface is done with PyQt6. The choice of the PyQt6 library was made because of its built-in support for multithreaded worker patterns, which is important for keeping the user interface responsive when time-consuming inference operations are performed. The construction is done with VS Code \cite{baumer2015visual} and the versioning is done with GitHub. \cite{dabbish2012social}.

\subsubsection{Hardware Feasibility}
Training deep models with long inference operations on complete hyperspectral cubes requires high computational resources. Google Colab \cite{bisong2019google}, provides the needed graphic processing units and eliminates the need for the computational resources for the project. The graphic user interface and integration design are done on the project team’s personal macOS systems.

\subsection{Legal Feasibility}
All libraries in the stack of PyQt6, Python, PyTorch, and OpenCV are fully open-sourced. The three benchmark datasets in use are legally permitted for academic use. There are no licensing or patent barriers.

\subsection{Economic Feasibility}
The entire tool chain is free and open-source with Google Colab offering the graphic processing units at no cost. The project makes no additional capital investments with the teams’ personal systems.

\subsection{Workforce and Time Feasibility}
The project team of Tan Erciyas and Kerem Baydar has the required skills to execute the entire project.

\begin{figure}[htbp]
	\centering

	\begin{ganttchart}[
		hgrid,
		vgrid,
		x unit=2cm,
		y unit title=0.6cm,
		y unit chart=1.2cm,
		title/.style={fill=gray!20, draw=black, font=\small\bfseries},
		bar/.style={fill=blue!70},
		bar height=0.4,
		bar label node/.append style={align=right, text width=4.8cm}
		]{1}{5}

		\gantttitle{2026}{5} \\
		\gantttitle{Feb}{1}
		\gantttitle{Mar}{1}
		\gantttitle{Apr}{1}
		\gantttitle{May}{1}
		\gantttitle{Jun}{1} \\

		\ganttbar{Literature Review \& Data Prep}{1}{1} \\
		\ganttbar{Preprocessing (PCA \& Filtering)}{2}{2} \\
		\ganttbar{Model Development (CNN \& ViT)}{2}{3} \\
		\ganttbar{Curriculum Learning Integration}{3}{4} \\
		\ganttbar{GUI \& Multithreading\\ Implementation}{4}{5} \\
		\ganttbar{Final Testing \& Documentation}{5}{5}
	\end{ganttchart}
	\caption{Project Gantt Chart (Feb 2026 - June 2026)}
\end{figure}
